% Este archivo ha sido generado automáticamente por generatabla.py
\section{Sprint 1: Unity, Javascript y MediaPipe}
\label{anexo3:sec:sprint1}

Explicación del sprint...

\sprinttable{
    nombre={Unity, Javascript y MediaPipe},
    id={1},
    tabla={
        Detector de pose con MediaPipe en Javascript       & 4 \\
        Comunicar el editor de Unity y JavaScript          & 3 \\
        Comunicar el export de WebGL de Unity y JavaScript & 3 \\
    },
    total={10},
    anadidos={},
    restantes={},
    completados={},
}

\diagrama{
    nombre={Burndown del sprint: Unity, Javascript y MediaPipe},
    archivo={anexo3/sprints/sprint1.pdf},
    etiqueta={anexo3:fig:burndown1},
    width=\textwidth,
    forzarPosicion={si pls},
}


\section{Sprint 2: Segundo Sprint}
\label{anexo3:sec:sprint2}

Explicación del sprint...

\sprinttable{
    nombre={Segundo Sprint},
    id={2},
    tabla={
        Use the position in image for estimating hip position & 2 \\
        Smooth out point position                             & 3 \\
        Change NN model depending on device                   & 1 \\
        Better mediapipe script                               & 2 \\
        Allow moving around                                   & 3 \\
        Try NextJS static export                              & 1 \\
    },
    total={12},
    anadidos={},
    restantes={},
    completados={},
}

\diagrama{
    nombre={Burndown del sprint: Segundo Sprint},
    archivo={anexo3/sprints/sprint2.pdf},
    etiqueta={anexo3:fig:burndown2},
    width=\textwidth,
    forzarPosicion={si pls},
}


\section{Sprint 3: Globos, objetivos y alternativas}
\label{anexo3:sec:sprint3}

Explicación del sprint...

\sprinttable{
    nombre={Globos, objetivos y alternativas},
    id={3},
    tabla={
        Obtener los objetivos del sistema                         & 2 \\
        Minijuego de globos                                       & 6 \\
        Enviar datos desde Unity a Javascript                     & 3 \\
        Situar la cámara en un overlay                            & 3 \\
        Probar OpenPose                                           & 3 \\
        Probar a utilizar solo los landmarks 2D para la detección & 3 \\
    },
    total={20},
    anadidos={},
    restantes={},
    completados={},
}

\diagrama{
    nombre={Burndown del sprint: Globos, objetivos y alternativas},
    archivo={anexo3/sprints/sprint3.pdf},
    etiqueta={anexo3:fig:burndown3},
    width=\textwidth,
    forzarPosicion={si pls},
}


\section{Sprint 4: Errores, limpieza y comienzo del backend}
\label{anexo3:sec:sprint4}

Explicación del sprint...

\sprinttable{
    nombre={Errores, limpieza y comienzo del backend},
    id={4},
    tabla={
        Crear base de datos de usuarios                & 5 \\
        El sistema no detecta las posturas al exportar & 1 \\
        Permitir saltar                                & 2 \\
        Combinar los dos cuerpos y limpiar el código   & 5 \\
        Arreglar el minijuego para presentar           & 3 \\
    },
    total={16},
    anadidos={},
    restantes={5},
    completados={11},
}

\diagrama{
    nombre={Burndown del sprint: Errores, limpieza y comienzo del backend},
    archivo={anexo3/sprints/sprint4.pdf},
    etiqueta={anexo3:fig:burndown4},
    width=\textwidth,
    forzarPosicion={si pls},
}


\section{Sprint 5: Backend simple}
\label{anexo3:sec:sprint5}

Explicación del sprint...

\sprinttable{
    nombre={Backend simple},
    id={5},
    tabla={
        Crear base de datos de usuarios & 5 \\
        Backend de usuario simple       & 6 \\
        Gestión de usuarios             & 10 \\
    },
    total={21},
    anadidos={},
    restantes={},
    completados={},
}

\diagrama{
    nombre={Burndown del sprint: Backend simple},
    archivo={anexo3/sprints/sprint5.pdf},
    etiqueta={anexo3:fig:burndown5},
    width=\textwidth,
    forzarPosicion={si pls},
}


\section{Sprint 6: Vite, Mantine y backend completo}
\label{anexo3:sec:sprint6}

Explicación del sprint...

\sprinttable{
    nombre={Vite, Mantine y backend completo},
    id={6},
    tabla={
        Crear base de datos completa                         & 2 \\
        Mejores excepciones en la API                        & 2 \\
        Todos los endpoints del backend                      & 3 \\
        Verificación de correo electrónico                   & 1 \\
        A Vite y Mantine                                     & 5 \\
        Montar pytest                                        & 2 \\
        Crear tests para las actividades                     & 3 \\
        Mover la API a un contenedor de Docker               & 1 \\
        Hacer que el DAO devuelva excepciones en vez de None & 1 \\
    },
    total={20},
    anadidos={7},
    restantes={},
    completados={},
}

\diagrama{
    nombre={Burndown del sprint: Vite, Mantine y backend completo},
    archivo={anexo3/sprints/sprint6.pdf},
    etiqueta={anexo3:fig:burndown6},
    width=\textwidth,
    forzarPosicion={si pls},
}


\section{Sprint 7: Frontend!}
\label{anexo3:sec:sprint7}

Explicación del sprint...

\sprinttable{
    nombre={Frontend!},
    id={7},
    tabla={
        Utilizar openapi-ts para generar el cliente de la API           & 2 \\
        Mover el frontend a un contendor de Docker                      & 2 \\
        Página de registro                                              & 4 \\
        Crear script para llenar la base de datos con valores de prueba & 2 \\
        Rutas protegidas                                                & 3 \\
        Página personal                                                 & 5 \\
        Página de estadísticas                                          & 8 \\
    },
    total={26},
    anadidos={},
    restantes={},
    completados={},
}

\diagrama{
    nombre={Burndown del sprint: Frontend!},
    archivo={anexo3/sprints/sprint7.pdf},
    etiqueta={anexo3:fig:burndown7},
    width=\textwidth,
    forzarPosicion={si pls},
}


\section{Sprint 8: Un poco de Unity}
\label{anexo3:sec:sprint8}

Explicación del sprint...

\sprinttable{
    nombre={Un poco de Unity},
    id={8},
    tabla={
        Limpiar el código del cuerpo                                      & 3 \\
        Enviar tamaño del video al juego de unity además de los landmarks & 1 \\
        Calcular los bounds del cuerpo                                    & 2 \\
        Detectar distancia a la cámara                                    & 4 \\
        Escena de calibración                                             & 8 \\
        Crear actividades desde Unity                                     & 2 \\
        Permitir tocar botones con el cuerpo en el juego de Unity         & 4 \\
    },
    total={24},
    anadidos={},
    restantes={},
    completados={},
}

\diagrama{
    nombre={Burndown del sprint: Un poco de Unity},
    archivo={anexo3/sprints/sprint8.pdf},
    etiqueta={anexo3:fig:burndown8},
    width=\textwidth,
    forzarPosicion={si pls},
}


\section{Sprint 9: Más Unity}
\label{anexo3:sec:sprint9}

Explicación del sprint...

\sprinttable{
    nombre={Más Unity},
    id={9},
    tabla={
        Menú principal en unity                                                                          & 3 \\
        Utilizar Movenet en vez de Mediapipe                                                             & 4 \\
        Utilizar solo los landmarks en dos dimensiones                                                   & 5 \\
        Cambiar iconos del frontend para que no tarde tanto en cargar                                    & 1 \\
        Acabar versión básica del juego de globos                                                        & 5 \\
        Arreglar los bounds del cuerpo                                                                   & 2 \\
        Arreglar texto del juego de globos (transparencia atascada?)                                     & 2 \\
        Combinar JSConnector y GameController en uno solo                                                & 2 \\
        Hacer que RINLBody obtenga los datos del JSConnector cuando los necesite para mejorar la fluidez & 1 \\
    },
    total={25},
    anadidos={2},
    restantes={},
    completados={},
}

\diagrama{
    nombre={Burndown del sprint: Más Unity},
    archivo={anexo3/sprints/sprint9.pdf},
    etiqueta={anexo3:fig:burndown9},
    width=\textwidth,
    forzarPosicion={si pls},
}


\section{Sprint 10: Un poco de todo}
\label{anexo3:sec:sprint10}

Explicación del sprint...

\sprinttable{
    nombre={Un poco de todo},
    id={10},
    tabla={
        Añadir scope a los fixtures de los tests                                               & 1 \\
        Utilizar cookies en vez de localStorage para el token de acceso                        & 2 \\
        Modificar endpoint del usuario                                                         & 2 \\
        Cliente propio para las requests                                                       & 3 \\
        Ordenar y filtrar los datos de actividad y forma física en el backend                  & 1 \\
        Permitir añadir más formas                                                             & 2 \\
        No permitir realizar acciones hasta que se haya verificado el correo                   & 1 \\
        No permitir jugar al juego hasta que se haya añadido una forma física                  & 1 \\
        Modificar datos del usuario                                                            & 1 \\
        Ocultar la barra de navegación y eliminar márgenes (o usar pantalla completa) al jugar & 1 \\
        Utilizar notificaciones en más sitios                                                  & 1 \\
        Crear herramientas para notificaciones                                                 & 2 \\
        Mutaciones de react query y useAuth                                                    & 1 \\
        Añadir campo puntiación a las actividades                                              & 2 \\
        Al iniciar el juego, enviar la forma física del usuario para calcular la actividad     & 1 \\
        Probar animaciones 2D                                                                  & 4 \\
        Send activity data, complete activity and return to main menu in ballon game           & 2 \\
    },
    total={28},
    anadidos={},
    restantes={6},
    completados={22},
}

\diagrama{
    nombre={Burndown del sprint: Un poco de todo},
    archivo={anexo3/sprints/sprint10.pdf},
    etiqueta={anexo3:fig:burndown10},
    width=\textwidth,
    forzarPosicion={si pls},
}


\section{Sprint 11: Más unity}
\label{anexo3:sec:sprint11}

Explicación del sprint...

\sprinttable{
    nombre={Más unity},
    id={11},
    tabla={
        Probar animaciones 2D                                       & 4 \\
        Calcular puntos de actividad                                & 5 \\
        Cambiar la calibración para que sea un guión de SceneScript & 3 \\
        El timer de los globos se ve raro                           & 1 \\
        Herramientas para enviar mensajes entre Unity y JS          & 2 \\
        El cuerpo está torcido solo en producción                   & 3 \\
        Arreglar que la puntuación siempre sea 0                    & 2 \\
        Hacer que los globos no aparezcan sobre las manos           & 2 \\
        Permitir cambiar la dificultad de los globos                & 3 \\
        Rotar el cuerpo hace que se rompa                           & 2 \\
        Mover la mano de la muñeca al centro de la mano de verdad   & 3 \\
    },
    total={30},
    anadidos={},
    restantes={},
    completados={},
}

\diagrama{
    nombre={Burndown del sprint: Más unity},
    archivo={anexo3/sprints/sprint11.pdf},
    etiqueta={anexo3:fig:burndown11},
    width=\textwidth,
    forzarPosicion={si pls},
}


\section{Sprint 12: Aun más Unity}
\label{anexo3:sec:sprint12}

Explicación del sprint...

\sprinttable{
    nombre={Aun más Unity},
    id={12},
    tabla={
        Enviar datos de actividad al finalizar las actividades          & 2 \\
        Suavizar movimiento del jugador                                 & 2 \\
        Hacer que el jugador se mantenga mejor en el suelo              & 2 \\
        Detectar cuando el jugador se acerca o aleja mucho de la cámara & 4 \\
        Añadir progreso al hacer click                                  & 3 \\
        Permitir apuntar a los botones en vez de tocarlos               & 3 \\
        Centrar cuerpo con los bounds en vez de con la cámara           & 2 \\
        Detectar cuando el jugador sale de los bounds                   & 2 \\
    },
    total={20},
    anadidos={},
    restantes={},
    completados={},
}

\diagrama{
    nombre={Burndown del sprint: Aun más Unity},
    archivo={anexo3/sprints/sprint12.pdf},
    etiqueta={anexo3:fig:burndown12},
    width=\textwidth,
    forzarPosicion={si pls},
}


\section{Sprint 13: Cosas varias y pensar en el nuevo minijuego}
\label{anexo3:sec:sprint13}

Explicación del sprint...

\sprinttable{
    nombre={Cosas varias y pensar en el nuevo minijuego},
    id={13},
    tabla={
        Cachear las búsquedas de actividades                                        & 1 \\
        Enpoint para reenviar verificación de email                                 & 1 \\
        Enviar correo de verificación al modificar el email                         & 2 \\
        Form de los datos del usuario                                               & 2 \\
        Cambiar el tiempo de chaché de las queries de react-query                   & 1 \\
        Organizar scripts de Unity                                                  & 1 \\
        Mover los bounds del cuerpo al centrarlo                                    & 2 \\
        Añadir una nueva forma física debería invalidar el caché de las actividades & 1 \\
        Añadir un id único a las landmarks para evitar convertirlas varias veces    & 1 \\
        Pensar y probar cosas para otro minijuego                                   & 6 \\
    },
    total={18},
    anadidos={},
    restantes={},
    completados={},
}

\diagrama{
    nombre={Burndown del sprint: Cosas varias y pensar en el nuevo minijuego},
    archivo={anexo3/sprints/sprint13.pdf},
    etiqueta={anexo3:fig:burndown13},
    width=\textwidth,
    forzarPosicion={si pls},
}


\section{Sprint 14: Nuevo minijuego}
\label{anexo3:sec:sprint14}

Explicación del sprint...

\sprinttable{
    nombre={Nuevo minijuego},
    id={14},
    tabla={
        Sistema de alertas para unity                              & 4 \\
        Test para minijuego de esquivar                            & 8 \\
        Intro para minijuego de esquivar                           & 4 \\
        Escoger dificultad para esquivar                           & 4 \\
        Fin de minijuego de esquivar (terminar actividad y puntos) & 3 \\
        Los pies del juegador se unden?                            & 4 \\
    },
    total={27},
    anadidos={},
    restantes={3},
    completados={24},
}

\diagrama{
    nombre={Burndown del sprint: Nuevo minijuego},
    archivo={anexo3/sprints/sprint14.pdf},
    etiqueta={anexo3:fig:burndown14},
    width=\textwidth,
    forzarPosicion={si pls},
}


\section{Sprint 15: Cosas varias}
\label{anexo3:sec:sprint15}

Explicación del sprint...

\sprinttable{
    nombre={Cosas varias},
    id={15},
    tabla={
        Fin para el minijuego de esquivar                                         & 3 \\
        Arreglar problema de CORS al enviar correos                               & 3 \\
        Penalizar por tocar globos con la mano equivocada                         & 3 \\
        Forzar un logout al cambiar el correo está mal                            & 3 \\
        El sexo del usuario no se usa para nada, se puede eliminar                & 2 \\
        Utilizar los bounds de la cámara en vez de perdir al usuario que se mueva & 3 \\
        Utilizar object pools para los globos                                     & 2 \\
        Obstáculos visibles en el minijuego de esquivar                           & 4 \\
    },
    total={23},
    anadidos={},
    restantes={},
    completados={},
}

\diagrama{
    nombre={Burndown del sprint: Cosas varias},
    archivo={anexo3/sprints/sprint15.pdf},
    etiqueta={anexo3:fig:burndown15},
    width=\textwidth,
    forzarPosicion={si pls},
}


\section{Sprint 16: Cosas varias 2}
\label{anexo3:sec:sprint16}

Explicación del sprint...

\sprinttable{
    nombre={Cosas varias 2},
    id={16},
    tabla={
        La información es algo sensible, estaría bien encriptarla         & 2 \\
        Tests para las formas físicas                                     & 1 \\
        Tests para los usuarios                                           & 2 \\
        Mejora menú principal de Unity                                    & 3 \\
        Evitar que el mismo tipo de globo salga demasiadas veces seguidas & 2 \\
        Hacer un objeto del canvas del mismo tamaño que los bounds        & 4 \\
    },
    total={14},
    anadidos={},
    restantes={5},
    completados={9},
}

\diagrama{
    nombre={Burndown del sprint: Cosas varias 2},
    archivo={anexo3/sprints/sprint16.pdf},
    etiqueta={anexo3:fig:burndown16},
    width=\textwidth,
    forzarPosicion={si pls},
}


\section{Sprint 17: Progreso?}
\label{anexo3:sec:sprint17}

Explicación del sprint...

\sprinttable{
    nombre={Progreso?},
    id={17},
    tabla={
        Página principal del dashboard                                    & 5 \\
        Mejorar menú principal                                            & 3 \\
        Evitar que el mismo tipo de globo salga demasiadas veces seguidas & 2 \\
        Tirar excepciones HTTP en vez de devolver respuestas con códigos  & 1 \\
        Sonidos para globos                                               & 1 \\
        Error al entrar en pantalla completa                              & 3 \\
        Mostrar claramente cuando no hay estadísticas                     & 2 \\
    },
    total={17},
    anadidos={},
    restantes={9},
    completados={8},
}

\diagrama{
    nombre={Burndown del sprint: Progreso?},
    archivo={anexo3/sprints/sprint17.pdf},
    etiqueta={anexo3:fig:burndown17},
    width=\textwidth,
    forzarPosicion={si pls},
}


\section{Sprint 18: Toca TFG}
\label{anexo3:sec:sprint18}

Explicación del sprint...

\sprinttable{
    nombre={Toca TFG},
    id={18},
    tabla={
        Página principal del dashboard                                   & 5 \\
        Mejorar menú principal de Unity                                  & 3 \\
        Tirar excepciones HTTP en vez de devolver respuestas con códigos & 1 \\
        Aprender algo de Latex                                           & 5 \\
        Arreglar errores del backend                                     & 4 \\
    },
    total={18},
    anadidos={},
    restantes={},
    completados={},
}

\diagrama{
    nombre={Burndown del sprint: Toca TFG},
    archivo={anexo3/sprints/sprint18.pdf},
    etiqueta={anexo3:fig:burndown18},
    width=\textwidth,
    forzarPosicion={si pls},
}


\section{Sprint 19: Documentación y arreglar cosas}
\label{anexo3:sec:sprint19}

Explicación del sprint...

\sprinttable{
    nombre={Documentación y arreglar cosas},
    id={19},
    tabla={
        Documentar backend para sphinx                                   & 6 \\
        Requisitos no funcionales                                        & 3 \\
        Página de 404                                                    & 2 \\
        Arreglar el docker compose del backend                           & 3 \\
        Sphinx                                                           & 4 \\
        Permitir pulsar botones con las dos manos                        & 5 \\
        Tablas de objetivos                                              & 4 \\
        Cambiar todo a naranja (como strava)                             & 1 \\
        Arreglar el despliegue                                           & 1 \\
        Bibliografía y glosario en latex                                 & 2 \\
        Tablas de requisitos de información                              & 3 \\
        Cambiar las tablas de IRQ a las nuevas y cambiar la bibliografía & 2 \\
        Tablas de actores                                                & 3 \\
        Tablas de participantes                                          & 2 \\
        Arreglar portadas documentos                                     & 2 \\
        Hacer comando para las tablas de caso de uso                     & 2 \\
        Arreglar referencias para tablas                                 & 2 \\
        Acabar la introducción                                           & 3 \\
    },
    total={50},
    anadidos={18},
    restantes={},
    completados={},
}

\diagrama{
    nombre={Burndown del sprint: Documentación y arreglar cosas},
    archivo={anexo3/sprints/sprint19.pdf},
    etiqueta={anexo3:fig:burndown19},
    width=\textwidth,
    forzarPosicion={si pls},
}


